%% kapa: exam=1 round=31 year=2020 period=1 form=A subject=civil_law subject_ko=민법 seq=19 number=19
%% source: Q-Net 감정평가사 기출문제 (cst00309) artlSeq=5207219 file=extracted/1008/5207219/A형/A형/2020년도 제31회 감정평가사 1차 시험_1교시_A형.pdf
%% conversion: kordoc 4.24.51 -> tools/gen1.py (auto); review=proofread-sonnet; flags=none
\begin{question}{19}
소멸시효와 등기 없는 취득시효에 관한 설명으로 옳은 것은? (다툼이 있으면 판례에 따름)
\begin{choices}
  \item 취득시효기간 동안 계속하여 등기명의인이 동일한 때에도 반드시 점유를 개시한 때를 기산점으로 하여야 한다.
  \item 점유자가 전(前)점유자의 점유를 아울러 주장할 때에는 그 점유의 개시시기를 어느 점유자의 점유기간 중 임의의 시점으로 선택할 수 있다.
  \item 채권의 소멸시효가 완성하면 그 채무자의 다른 채권자는 직접 그 완성을 원용할 수 있다.
  \item 압류 또는 가압류는 소멸시효와 취득시효의 중단사유이다.
  \item 취득시효의 중단사유는 종래의 점유상태의 계속을 파괴하는 것으로 인정될 수 있는 것이어야 한다.
\end{choices}
\end{question}
\answer{5}
